%% LyX 2.4.4 created this file.  For more info, see https://www.lyx.org/.
%% Do not edit unless you really know what you are doing.
\documentclass[english]{article}
\usepackage[T1]{fontenc}
\usepackage[latin9]{inputenc}
\usepackage{cancel}
\usepackage{geometry}
\geometry{verbose,tmargin=1in,bmargin=1in,lmargin=1in,rmargin=1in}
\PassOptionsToPackage{normalem}{ulem}
\usepackage{ulem}

\makeatletter

%%%%%%%%%%%%%%%%%%%%%%%%%%%%%% LyX specific LaTeX commands.
%% Because html converters don't know tabularnewline
\providecommand{\tabularnewline}{\\}

%%%%%%%%%%%%%%%%%%%%%%%%%%%%%% User specified LaTeX commands.
\usepackage{hyperref}
\usepackage[usenames,dvipsnames]{xcolor} %adds additional color options
\hypersetup{colorlinks=true, urlcolor=Blue}

\usepackage{sectsty} %%one way to change section header font sizes.
\sectionfont{\large}

\usepackage{titlesec}
\titlespacing*{\section}{0pt}{10pt}{0pt}

\usepackage{multicol}

\usepackage{tikz}
\usetikzlibrary{plotmarks}
\usepackage{pgfplots}
\pgfplotsset{compat=newest}

\usepackage{calc}
\usepackage[nomessages]{fp}

\usepackage{advdate}    % Advancing/saving dates
\usepackage{datetime}   % Dates formatting
\usepackage{datenumber} % Counters for dates
%\longdate %\usdate %\DM %\MY
\newdateformat{MD}{\monthname[\THEMONTH] \ordinal{DAY}}
% Added by lyx2lyx
\setlength{\parskip}{\medskipamount}
\setlength{\parindent}{0pt}

\makeatother

\usepackage{babel}
\begin{document}
\begin{center}
{\Large\textbf{Johnson C. Smith University}}\textbf{}\\
\textbf{Finite Math~~|~~MTH 132~~|~~Section C~~|~~Fall
2012}\\
\textbf{Tuesday 9:30-10:45~~|~~MAIN SHA 108}\\
\textbf{\rule[0.25ex]{1\columnwidth}{1pt}}
\par\end{center}

\vspace{-15pt}

\begin{center}
\emph{\textquotedblleft Every act of conscious learning requires the
willingness to suffer an injury to one's self-esteem. That is why
young children, before they are aware of their own self-importance,
learn so easily; and why older persons, especially if vain or important,
cannot learn at all.\textquotedblright{}} \\
--Thomas Szasz, author, professor of psychiatry (1920- 2012)\vspace{-.25cm}
\par\end{center}

\section*{Instructor Information}

\textbf{Name: }Dr.\,James Ashe\\
\textbf{Office location: }307 Davis Hall (back)\\
\textbf{Email: }\href{mailto:jashe@jcsu.edu}{jashe@jcsu.edu}\\
\textbf{Office phone: }(704)\,378-1037\\
\textbf{Office hours: }TR 8:30-9:45 and 10:45-12; TR 4:15-7:45 (by
appointment)

\section*{Course Description}

This is a 3 credit hour\textbf{\emph{ }}course. Elements of finite
mathematical systems for liberal arts and education students. Topics
include real numbers, linear equations and straight ines, systems
of linear equations and inequalities, matrix algebra, sets and counting,
concepts of probability and statistics, and mathematics of finance.
The course relies heavily on computers and graphing calculators to
develop intuition, make estimates, verify results, and check reasonableness
of answers. \textbf{Prerequisite: MTH 131}. 

This is a \emph{hybrid course}. Hybrid courses meet approximatelt
49\% online and 51\% face-to-face. Online materials will be dissemated
through \emph{Jenzabar} \href{https://my.jcsu.edu/ics}{https://my.jcsu.edu/ics}.

Instructional methods will primarily consist of lectures (via online
videos), in-class discussions, and projects. However other methods
may be adopted according to the particular needs of the class.

\section*{Course Materials }

\textbf{Textbook/MyMathLab: }A MyMathLab access code is required.
New textbooks and Ebooks usually come packaged with a MyMathLab code.
\uline{Course ID: ashe58384}.\\
Register and purchase a code here: \href{http://www.pearsonmylabandmastering.com/northamerica/mymathlab/students/get-registered/index.html}{http://www.pearsonmylabandmastering.com/northamerica/mymathlab/students/get-registered/index.html}\\
An access code can also be purchased here:\\
 \href{http://www.amazon.com/MyMathLab-Student-Hall-Pearson-Education/dp/032119991X}{http://www.amazon.com/MyMathLab-Student-Hall-Pearson-Education/dp/032119991X}

\textbf{Calculator: }TI-83, TI-83+, and TI-84+ are allowed and supported.
Other calculators can be used with Instructor permission. Students
should bring their calculators to class.

\section*{Grades}

Grades will be taken from homework/quizzes, exams, a project, a common
Midterm, and common Final as follows:\textbf{}\\
\textbf{Evaluation: }HW/Quizzes $=20\%$, Exams 1, 2, 3, 4, and Midterm
$=12\%$ each; Final Exam $=20\%$.

\textbf{Grading scale: $\mbox{A}\geq90>\mbox{B}\geq80>\mbox{C}\geq70>\mbox{D}\geq60>\mbox{F}$}

\textbf{Quizzes: }HW/Quizzes will consist of any graded assignment
which is not an exam. This includes a \emph{project based learning
module}, traditional quizzes, MyMathLab assignments, take home assignments,
and in-class projects. The two lowest quiz grade will be dropped.
\textbf{}\\
\textbf{Final Exam and Midterm: }Both the final and midterm exam will
be common comprehensive exams given in-class; see the exam schedule
\href{http://www.jcsu.edu/happenings/exam-schedule}{jcsu.edu/happenings/exam-schedule}.
\emph{If beneficial, the final exam score will replace the lowest
exam score.}

\section*{Class Policies}

JCSU policies regarding attendance, academic integrity, grades, and
change of enrollment will be enforced. It is the student\textquoteright s
responsibility to be aware of these policies and of related deadlines
as well as any announcements or syllabus adjustments (written or oral)
made in class. Please see the student handbook for details, \href{http://www.jcsu.edu/docs/handbook.pdf}{jcsu.edu/docs/handbook.pdf}.

\textbf{Attendance: }Johnson C. Smith University has no official attendance
policy. However, because attendance in classes is a vital part of
the educational process, students are encouraged to attend classes
regularly and promptly. 

Students who may miss classes while representing the University in
an official capacity are exempt from regulations governing absences.
However, absence from class for official University business does
not relieve the student from responsibility for any class assignments
that may be missed during the period of absence.

\textbf{Missed quizzes: }There will be no make-ups for in-class quizzes
or assignments unless the class was missed for a University sponsored
event, and advance notice of the absence was given by the student.
In this case the assignment will be excused or a make-up will be given
depending on circumstances. Late take-home assignments will have 25\%
deducted for each day it is late. \\
\textbf{Missed exams: }If a student misses an exam for a verifiable
excuse, they should send me an e-mail \emph{as soon as possible }so
that a make-up can be arranged. Make-up exams may be significantly
more difficult. If you will be gone for any reason, it will be \emph{your
}responsibility to inform me at least a week in advance to schedule
to take the exam early. 

\textbf{Student support services: }If you need course adaptations
or accommodations for a documented disability please contact the Office
of Disability Services. \\
\href{http://www.jcsu.edu/directory/federal-requirements/general-institutional-information/jcsu-disability-services}{jcsu.edu/directory/federal-requirements/general-institutional-information/jcsu-disability-services},
\\
phone: (704)\,398-1116.

The Office of Counseling Services provides confidential intakes and
assessments, personal counseling and referrals to appropriate resources
to all enrolled students on a walk-in or appointment basis. \\
\href{http://www.jcsu.edu/admissions/future_students/meet_your_counselor}{jcsu.edu/admissions/future\_students/meet\_your\_counselor},
phone (704)\,378-1044.

\begin{multicols}{2}

\section*{Course Content}

\textbf{Chapter 1: Linear Functions} \\
1.1 Slopes and Equations of Lines \\
1.2 Linear Functions \& Applications\\
1.3 Least Squares Line

\textbf{Chapter 2: Systems of Linear Equations}\\
2.1 The Echelon Method\\
2.2 The Gauss-Jordan Method\\
2.3 Adding \& Subtracting Matrices\\
2.4 Matrix Multiplication\\
2.5 Matrix Inverses

\textbf{Chapter 5: Math of Finances}\\
5.1 Simple and Compound Interest\\
5.2 Future Value of Annuity\\
5.3 Present Value \& Amortization of Annuity

\textbf{Chapter 6: Logic}\\
6.1 Statements \\
6.2 Truth Tables \& Equivalent Statements\\
6.4 More on the Conditional

\textbf{Chapter 7: Sets and Probability}\\
7.1 Sets\\
7.2 Intoduction to Probability\\
7.3 Basics of Probability

%%\vfill
\columnbreak 

\section*{Important Dates}

\noindent\begin{minipage}[t]{1\columnwidth}%
\begin{tabular}{rcl}
\MD\formatdate{17}{1}{2014} &  & Last day to Add/Drop \tabularnewline
\MD\formatdate{28}{1}{2014} &  & Exam 1\tabularnewline
\MD\formatdate{18}{2}{2014} &  & Exam 2\tabularnewline
\MD\formatdate{4}{3}{2014} &  & Midterm Exam\tabularnewline
\MD\formatdate{1}{4}{2014} &  & Exam 3\tabularnewline
\MD\formatdate{15}{4}{2014} &  & Exam 4\tabularnewline
\MD\formatdate{30}{4}{2014} &  & Common Final (1-3pm)\tabularnewline
 &  & \tabularnewline
\end{tabular}%
\end{minipage}

The above exam dates are subject to change. JCSU's academic calendar
can be found here: \href{http://www.jcsu.edu/happenings/academic-calendar}{jcsu.edu/happenings/academic-calendar},
and the exam schedule can be found here: \href{http://www.jcsu.edu/happenings/exam-schedule/spring-exam}{http://www.jcsu.edu/happenings/exam-schedule/spring-exam}\textbf{Disclaimer:
}While I do not foresee any significant changes to the items in this
syllabus during the semester, I reserve the right to make appropriate
changes. Students will be notified of any such changes in writing.

\end{multicols}
\end{document}
